\section{Introduction}
\label{sec:introduction}

Test-time training can commit a parameter update before downstream
feedback tells us whether the update helped. When that feedback calls for
going back, the system needs both the earlier parameters and the optimizer
buffers. A full checkpoint makes restoration a copy operation, while a
reversible update can reconstruct the state from a tape of event
identifiers by recomputing gradients. The tradeoff depends on how much
state would need to be saved and how often the policy would need it back.

We study that tradeoff in analytic and small neural probes, then in a
controlled branching task with a causal Transformer. The experiments
separate the accuracy of reconstructing an optimizer state from the
effect of using it. They also test whether geometric priors preserve
prior-task accuracy across a continual stream.

In real arithmetic, we can run the same positive-step position-Verlet
update backward by reflecting momentum and reversing the batch order.
Its checked fixed-point counterpart restores integer endpoints bit for bit,
using one fresh gradient per step, as long as it remains in range.
Exact inversion preserves the rounded map, including kicks that round to
zero. Native floating-point replay has a different limitation. On the long
Transformer tape, its error need not grow according to the pre-specified
exponential model.

Feedback timing determines when reconstruction helps. A decision made
only after a window of speculative updates can leave earlier bad updates
in place, even if the policy stops at the alarm. When a warning is
available sooner, an exploratory comparison finds that skipping nearly
matches rewind in accuracy on the same tapes. The retention experiments
show a separate limit. Dissipating kinetic energy or contracting
displacements toward an anchor does not by itself preserve an old
task's accuracy.

In full-model search, replay lowers peak device allocation compared with
resident dense checkpoints. Pinned-host SGD is faster at the
matched deep-search solve rate, and an exploratory head-only control
shows how much the amount of trainable state matters. We report those
measurements separately from an analytical model of storage at larger
scales. The project repository is \RepoLink{}. Appendix~\ref{app:reproducibility}
identifies the build README and the archived revisions used to reproduce
the results.
